\subsection{分次模的分次张量积}

设 \(\Delta\) 是一个交换幺半群，其单位元记作 0，A 是一个类型为 \(\Delta\) 的分次环，M（相应地，N）是一个类型为 \(\Delta\) 的分次右（相应地，左）A-模。设 \((\mathbf{A}_{\lambda})\) （相应地， \((\mathbf{M}_{\lambda}), (\mathbf{N}_{\lambda}))\) ）是 A（相应地，M，N）的分次；Z-模 M 和 N 的张量积 \(\mathbf{M} \otimes_{\mathbf{Z}} \mathbf{N}\) 是诸 \(\mathbf{M}_{\lambda} \otimes_{\mathbf{Z}} \mathbf{N}_{\mu}\) 的直和（§3, 第 7 号, 命题 7），因而后者构成这个 Z-模上类型为 \(\Delta, \Delta\) 的一个双分次。考虑 M \(\otimes_{\mathbf{Z}} \mathbf{N}\) 上与这个双分次相关联的类型为 \(\Delta\) 的全分次（第 1 号，例 4）；它由子 Z-模 \(\mathrm{P}_{\lambda} = \sum_{\mu + v = \lambda} (\mathrm{M}_{\mu} \otimes_{\mathbf{z}} \mathrm{N}_{v})\) 组成。已知 Z-模 M \(\otimes_{\mathbf{A}} \mathbf{N}\) 是 M \(\otimes_{\mathbf{Z}} \mathbf{N}\) 关于由元素 \((xa) \otimes y - x \otimes (ay)\) （其中 \(x \in \mathbf{M}, y \in \mathbf{N}\) 和 \(a \in \mathbf{A}\) ，§3, 第 1 号）生成的子 Z-模 Q 的商；如果对所有 \(\lambda \in \Delta\) , \(x_{\lambda}, y_{\lambda}, a_{\lambda}\) 分别是次数为 \(\lambda\) 的齐次分量（ \(x, y, a\) 的），则显然 \((xa) \otimes y - x \otimes (zy)\) 是齐次元素 \((x_{\lambda}a_{v}) \otimes y_{\mu} - x_{\lambda} \otimes (a_{v}y_{\mu})\) 之和，换言之，Q 是 M \(\otimes_{\mathbf{Z}} \mathbf{N}\) 的一个分次子 Z-模（第 3 号, 命题 2），且商

\[
\mathbf {M} \otimes_ {\mathbf {A}} \mathbf {N} = (\mathbf {M} \otimes_ {\mathbf {Z}} \mathbf {N}) / \mathbf {Q}
\]

因此典范地具有一个类型为 \(\Delta\) 的分次 Z-模结构（第 3 号）。此外（第 3 号, 命题 5），A 的中心 C 是 A 的一个分次子环；我们刚在 M \(\otimes_{\mathbf{A}} \mathbf{N}\) 上定义的这个分次与其在分次环 C 上的模结构相容。因为 M \(\otimes_{\mathbf{Z}} \mathbf{N}\) 典范地具有两个 C-模结构，对于它们分别有 \(c(x \otimes y) = (xc) \otimes y\) 和 \((x \otimes y)c = x \otimes (cy)\) （对 \(x \in \mathbf{M}, y \in \mathbf{N}, c \in \mathbf{C}\) ）（ \(\S 3\) , 第 3 号）；若 \(x \in \mathbf{M}_{\lambda}, y \in \mathbf{N}_{\mu}, c \in \mathbf{C} \cap \mathbf{A}_{\nu}\) ，则两个元素 \(c(x \otimes y)\) 和 \((x \otimes y)c\) 属于 (M \(\otimes_{\mathbf{Z}} \mathbf{N})_{\lambda + \mu + \nu}\) ，且它们的差属于 Q，因此它们在 M \(\otimes_{\mathbf{A}} \mathbf{N}\) 中的共同像属于 (M \(\otimes_{\mathbf{A}} \mathbf{N})_{\lambda + \mu + \nu}\) ，这便确立了我们的断言。当我们把 M \(\otimes_{\mathbf{A}} \mathbf{N}\) 作为分次 C-模谈论时，除非另有说明，总是指如此定义的结构。注意 (M \(\otimes_{\mathbf{A}} \mathbf{N})_{\lambda}\) 可以定义为 M \(\otimes_{\mathbf{A}} \mathbf{N}\) 的由诸 \(x_{\mu} \otimes y_{\nu}\) （其中 \(x_{\mu} \in \mathbf{M}_{\mu}, y_{\nu} \in \mathbf{N}_{\nu}\) 且 \(\mu + \nu = \lambda\) ）生成的加法子群。

设 \(\mathbf{M}'\) （相应地， \(\mathbf{N}'\) ）是另一个分次右（相应地，左）A-模，且 \(u: \mathbf{M} \to \mathbf{M}', v: \mathbf{N} \to \mathbf{N}'\) 是次数分别为 \(\alpha\) 和 \(\beta\) 的分次同态。那么由上述注记立即得出 \(u \otimes v\) 是次数为 \(\alpha + \beta\) 的分次（C-模）同态。

当A是交换环时，在任意有限多个分次A-模的张量积上同样可以定义与A-模结构相容的分次结构；此外，显然诸如结合性同构 \((\mathbf{M} \otimes \mathbf{N}) \otimes \mathbf{P} \to \mathbf{M} \otimes (\mathbf{N} \otimes \mathbf{P})\) (§ 3，第 8 号，命题 8)是分次模的同构。

注. 当 A 具有平凡分次（第 1 号，例 1）时， \((\mathbf{M} \otimes_{\mathbf{A}} \mathbf{N})_{\lambda}\) 就是诸子 Y-模 \(\mathbf{M}_{\mu} \otimes_{\mathbf{A}} \mathbf{N}_{\nu}\) （ \(\mathbf{M} \otimes_{\mathbf{A}} \mathbf{N}\) 中的，满足 \(\mu + \nu = \lambda\) ）的直和。

设 M（相应地，N）是类型为 \(\Delta\) 的分次右（相应地，左）A-模，P 是类型为 \(\Delta\) 的分次 Z-模，并设 f 是 \(M \times N\) 到 P 中的一个 Z-双线性映射，它满足 §3, 第 1 号的条件 (1)，且此外使得

\[
f (x _ {\lambda}, y _ {\mu}) \in \mathrm{P} _ {\lambda + \mu} \quad \text { for } x _ {\lambda} \in \mathrm{M} _ {\lambda}, y _ {\mu} \in \mathrm{N} _ {\mu}, \lambda , \mu \text { in } \Delta .
\]

那么 \(f(x, y) = g(x \otimes y)\) ，其中 \(g: \mathbf{M} \otimes_{\mathbf{A}} \mathbf{N} \to \mathbf{P}\) 是一个 \(\mathbf{Z}\) -线性映射（§3, 第 1 号, 命题 1），且由上述条件可知 \(g\) 是次数为 0 的分次 \(\mathbf{Z}\) -模同态。

设 B 是另一个类型为 \(\Delta\) 的分次环， \(\rho: A \to B\) 是分次环的一个同态（第 2 号）；则 \(\rho_*(B_d)\) 是类型为 \(\Delta\) 的一个分次右 A-模。若 E 是类型为 \(\Delta\) 的分次左 A-模，且给 \(\rho^*(B_d) \otimes_A E\) 赋予上面定义的类型为 \(\Delta\) 的分次 Z-模结构，则典范的左 B-模结构（ \(\S\) 5, 第 1 号）与

\[
\mathbf {E} _ {(\mathbf {B})} = \rho^ {*} (\mathbf {E}) = \rho_ {*} (\mathbf {B} _ {d}) \otimes_ {\mathbf {A}} \mathbf {E}.
\]

的分次相容。这样得到的分次 B-模称为借助 \(\rho\) 把纯量环扩张到 B 而得到的；当我们把 \(\mathrm{E}_{(\mathrm{B})}\) 或 \(\rho^{*}(\mathrm{E})\) 作为分次 B-模谈论时，除非另有说明，总是指这一结构。

